\section{Objetivos}
Nuestro objetivo general es desarrollar y validar una metodología de creación un modelo de Automated Storage Tiering (AST) basado en aprendizaje automático,
para reducir la saturación de la capa activa de almacenamiento y minimizar el coste operativo, garantizando el cumplimiento de los acuerdos de nivel de servicio (SLA) en entornos de producción reales.


Los objetivos específicos son:
\begin{itemize}
    \item Diagnosticar el estado actual de las infraestructuras de almacenamiento analizadas (CTTI y Google Cluster Traces) y establecer una línea base cuantitativa de su ineficiencia porcentaje de capacidad inactiva, frecuencia de saturación de la capa activa y coste operativo asociado a la sobreestimación.
    \item Diseñar e implementar un modelo de ML para la clasificación automática de ficheros y bloques en capas de almacenamiento (hot, warm, cold), incorporando restricciones de SLA como condición operativa.
    \item Evaluar el impacto del modelo comparando las métricas (saturación de la capa activa, latencia de acceso, cumplimiento de SLA y coste operativo evitado por unidad de tiempo) frente a las políticas estáticas de HSM.
\end{itemize}

